% ============================================================================
% Chapter 9 — Conclusion and Future Work
% ============================================================================
\chapter{Conclusion and Future Work}
\label{ch:conclusion}

This chapter concludes the thesis by summarizing our primary research contributions and suggesting directions for future work.


% ---- 9.1 Summary of Contributions -------------------------------------------
\section{Summary of Contributions}
\label{sec:summary}

This thesis presented \versimux{}, a grounded multi-agent architecture for automated usability evaluation and code-level remediation of web interfaces. The system makes five primary contributions:

\begin{enumerate}
  \item \textbf{Grounded Persona Simulation:} We developed a multi-agent cognitive simulation model that interacts with web pages via browser automation. By coupling visual inputs and structural DOM states with grounding guards, we eliminated execution-time selector hallucinations, ensuring the agents remain synchronized with the application's actual runtime state.
  \item \textbf{Closed-Loop Remediation Pipeline:} We designed a multi-agent pipeline that automates the transition from diagnostic discovery to verified code repair. The system detects violations, synthesizes targeted CSS, HTML, and JS patches, applies the repairs to the target files, and re-runs the persona simulations to verify the effectiveness of the fixes.
  \item \textbf{Conflict-Aware Patch Synthesis:} We implemented a collaborative patching subsystem that resolves overlapping or contradictory remediation proposals. Using a multi-agent mediation and debate protocol, the conflict resolver node evaluates overlapping patches to synthesize a unified, conflict-free set of modifications.
  \item \textbf{Trace Integrity Verification:} We introduced an automated trace verification process that validates simulator outputs. This node identifies redundant interaction loops, filters out false positives, and ensures that the issues generated by the simulated personas correspond to verifiable interaction blocks.
  \item \textbf{Real-time Observability and Infrastructure:} We engineered a thread-safe infrastructure to support concurrent evaluation tasks. By integrating a singleton EventBus with SQLite WAL persistence and Server-Sent Events, the backend streams progress updates, logs, and screenshots to the Next.js dashboard, supporting real-time tracing and state recovery via the Last-Event-ID protocol.
\end{enumerate}


% ---- 9.2 Future Work --------------------------------------------------------
\section{Future Work}
\label{sec:future-work}

Several directions for future work emerge from this thesis:

\begin{description}
  \item[Multi-page Navigation:] Extending the persona agents to navigate across multiple stateful pages would enable the evaluation of multi-step user workflows (such as user onboarding or checkout funnels) and the detection of complex navigation blocks.
  \item[Visual Regression Testing:] Integrating automated screenshot-based visual regression checks into the verification layer would prevent generated CSS and HTML patches from introducing layout shifts or breaking visual designs.
  \item[Domain-Specific Fine-Tuned Models:] Fine-tuning lightweight open-source models (such as Llama 3.1 8B) on web accessibility datasets and DOM repair tasks could reduce API costs and inference latency, enabling offline execution.
  \item[CI/CD Workflow Integration:] Packaging \versimux{} as a GitHub Action or pre-commit hook would enable automated accessibility evaluations and proposed fixes to be generated directly inside developer pull requests.
  \item[Longitudinal Developer Adoption Studies:] Conducting longitudinal studies with web development teams would provide insights into how developers interact with automated reports and patches, helping to guide improvements in the user interface.
  \item[Direct JavaScript SPA Component Patching:] Extending the patching engine to write fixes directly into modern component trees (e.g., React, Vue, or Svelte component files) would eliminate the need for manual porting, integrating with modern frontend codebases.
\end{description}
